\documentclass[11pt]{article}
\usepackage[margin=1in]{geometry}
\usepackage{amsmath,amssymb,amsthm,hyperref}
\newtheorem{proposition}{Proposition}
\newtheorem{definition}{Definition}
\newcommand{\dd}{\,\mathrm d}
\title{Endogenous Task Automation: Initial Analytical Results}
\author{Working note for Oliver Pardo}
\date{22 November 2026}
\begin{document}
\maketitle
\noindent This note is separate from the existing manuscript. I establish a static
task allocation and conditional long-run implications. I do not claim dynamic
equilibrium existence, global research optimality, or novelty.

\section{Environment and the task allocation}
I retain the original household, capital accumulation, compute cost and, initially,
the bounded research technology. Write $H=AL$. Final production is
\begin{equation}
 Y=K^\alpha Z^{1-\alpha},\qquad
 Z=\left[\int_0^1 z(i)^{(\sigma-1)/\sigma}\dd i\right]^{\sigma/(\sigma-1)}.
 \label{eq:tasks-aggregate}
\end{equation}
The unit-elastic limit is $\log Z=\int_0^1\log z(i)\dd i$.
Within task $i$, human and AI services are perfect substitutes:
\begin{equation}
 z(i)=h(i)+e^{-\beta i}x(i),\qquad
 \int_0^1h(i)\dd i=H,\quad \int_0^1x(i)\dd i=X,\quad \beta>0.
 \label{eq:tasks-technology}
\end{equation}
Humans can move freely between tasks. The coefficient $e^{-\beta i}$ is fixed.
All tasks are technically automatable at every date. Research lowers their cost
through $X=BU$ rather than changing their technological feasibility.
I call the optimized composite $\mathcal Z(H,X)$.

\begin{proposition}[Efficient task assignment]
For $H,X>0$, the unique cutoff $m\in(0,1)$ satisfies
\begin{equation}
 \frac XH=\frac{D_\sigma(m)e^{\sigma\beta m}}{1-m},\qquad
 D_\sigma(m)=\int_0^m e^{-\beta(\sigma-1)i}\dd i.
 \label{eq:tasks-cutoff}
\end{equation}
AI supplies tasks $i<m$ and humans supply tasks $i>m$. At this allocation,
\begin{equation}
 \mathcal Z(H,X)=
 \left[(1-m)^{1/\sigma}H^{(\sigma-1)/\sigma}
       +D_\sigma(m)^{1/\sigma}X^{(\sigma-1)/\sigma}
 \right]^{\sigma/(\sigma-1)}.
 \label{eq:tasks-reduced}
\end{equation}
Equation~\eqref{eq:tasks-reduced} is conditional on the optimal cutoff.
It is not a CES with constant coefficients as $X/H$ changes.
\end{proposition}
\begin{proof}
At a candidate cutoff, allocate $h(i)=H/(1-m)$ on human tasks and
$x(i)=X e^{-\beta(\sigma-1)i}/D_\sigma(m)$ on AI tasks.
Their outputs coincide at the cutoff precisely when
equation~\eqref{eq:tasks-cutoff} holds. This equation has one solution because
the logarithmic derivative of its right-hand side is
\[
 D_\sigma'(m)/D_\sigma(m)+\sigma\beta+1/(1-m)>0,
\]
and its endpoint limits are zero and infinity.

To establish global efficiency, let $q(i)=\partial Z/\partial z(i)$ at the candidate.
On human tasks $q(i)$ is constant. On AI tasks $q(i)e^{-\beta i}$ is constant.
The common human return divided by the common AI return is $e^{\beta m}$.
Thus neither factor can raise its return by moving to another task, including
tasks currently supplied by the other factor. The CES is concave for every
$\sigma>0$. Its supporting tangent therefore bounds the output of every
alternative feasible allocation by the output of this candidate.
Substituting the candidate allocations into~\eqref{eq:tasks-aggregate} yields
equation~\eqref{eq:tasks-reduced}.
\end{proof}

At $\sigma=1$, the reduced expression is instead written as
\begin{equation}
 \log\mathcal Z=(1-m)\log\frac{H}{1-m}
     +m\log\frac Xm-\frac{\beta m^2}{2},\qquad
 \frac XH=\frac{m e^{\beta m}}{1-m}.
 \label{eq:tasks-unit}
\end{equation}

\section{Factor prices and adoption}
Final-good firms are competitive and pay $w$ per worker and $p_X$ per AI-service unit.
The final good is the numeraire. Their conditions give
\begin{equation}
 w=A(1-\alpha)(K/\mathcal Z)^\alpha\mathcal Z_H,\qquad
 p_X=(1-\alpha)(K/\mathcal Z)^\alpha\mathcal Z_X.
 \label{eq:tasks-prices}
\end{equation}
\begin{proposition}[The marginal task determines the wage--AI-price ratio]
At any allocation with positive AI use,
\begin{equation}
 \boxed{\frac{w}{A p_X}=e^{\beta m}.}
 \label{eq:tasks-wage-boundary}
\end{equation}
\end{proposition}
\begin{proof}
The ratio of the marginal products $\mathcal Z_H/\mathcal Z_X$ equals
the ratio of the task returns in the preceding proof, $e^{\beta m}$.
Taking the ratio of equations~\eqref{eq:tasks-prices} yields the result.
Equivalently, the human cost per task unit $w/A$ equals the AI cost
$p_Xe^{\beta m}$ at the cutoff.
\end{proof}

The inverse demand $P(X;K,H)$ is the second expression
in~\eqref{eq:tasks-prices}, evaluated at the efficient allocation.
The frontier responds to a seller's quantity, rather than remaining fixed
when I calculate the inverse-demand elasticity.

\begin{proposition}[Adoption threshold and the static monopoly restriction]
At given $K,H,B$, competitive supply at $p_X=1/B$ has positive adoption if and only if
\begin{equation}
 B(1-\alpha)(K/H)^\alpha>1.
 \label{eq:tasks-adoption-threshold}
\end{equation}
The developer's static problem
$\max_{X\geq0}\{P(X;K,H)X-X/B\}$ has a finite global maximizer.
It is positive under the same strict inequality. Every positive global
maximizer supplies fewer AI services and automates fewer tasks than competition
at the same $K,H,B$.
\end{proposition}
\begin{proof}
At $X=0$, $\mathcal Z(H,0)=H$ and its right derivative with respect to $X$ is one.
Thus $P(0;K,H)=(1-\alpha)(K/H)^\alpha$. The optimized composite is
increasing and concave in $X$. Since $K$ is fixed and $1-\alpha<1$,
the inverse demand is strictly decreasing on its positive domain.
It tends to zero at infinity. Competitive demand is positive exactly when
$P(0)>1/B$.

Monopoly operating profit is zero at $X=0$. If $P(0)\leq1/B$, every
positive $X$ yields negative profit. If $P(0)>1/B$, a sufficiently small
positive $X$ yields positive profit. At infinity, revenue is of order
$X^{1-\alpha}$ whereas cost is linear, so profit tends to minus infinity.
A positive global maximizer therefore exists and is interior. Its condition
$P+XP'=1/B$ implies $P>1/B$. Strictly decreasing demand places this quantity
below the competitive quantity. The cutoff increases strictly with $X/H$.
This argument does not establish uniqueness of the monopoly maximizer.
\end{proof}

\section{A conditional long-run result}
\begin{proposition}[Automation can remove the task bottleneck at any elasticity]
For every fixed $\sigma>0$, as $X/H\to\infty$,
\begin{equation}
 m\to1,\qquad \frac{\mathcal Z(H,X)}X\to c_\sigma,\qquad
 \mathcal Z_X\to c_\sigma,\qquad
 c_\sigma=\left[\int_0^1e^{-\beta(\sigma-1)i}\dd i\right]^{1/(\sigma-1)}.
 \label{eq:tasks-linear-limit}
\end{equation}
At $\sigma=1$, $c_1=e^{-\beta/2}$.
\end{proposition}
\begin{proof}
Equation~\eqref{eq:tasks-cutoff} implies $m\to1$.
For $\sigma\ne1$, set $S=1-m+D_\sigma(m)e^{\beta(\sigma-1)m}$.
Direct substitution in~\eqref{eq:tasks-reduced}, followed by differentiation
using the envelope theorem, gives
\[
 \frac{\mathcal Z}H=\frac{S^{\sigma/(\sigma-1)}}{1-m},\qquad
 \mathcal Z_H=S^{1/(\sigma-1)},\qquad
 \mathcal Z_X=e^{-\beta m}S^{1/(\sigma-1)}.
\]
Divide the first expression by~\eqref{eq:tasks-cutoff} and take $m\to1$.
For $\sigma=1$, differentiate~\eqref{eq:tasks-unit} at the optimum to obtain
$\mathcal Z_H=e^{\beta m^2/2}$ and
$\mathcal Z_X=e^{\beta m^2/2-\beta m}$, giving the same limits.
\end{proof}

This is a property of a bounded relative-difficulty distribution with no
permanently human-only tasks. It is not a claim that complementarity never
limits growth in other task models.

\begin{proposition}[A bound on long-run average wage growth]
On any path with positive AI use at all sufficiently late dates,
\begin{equation}
 \log\frac{w_t}{w_{t_0}}=\gamma(t-t_0)
 +\log\frac{p_{X,t}}{p_{X,t_0}}+\beta(m_t-m_{t_0}).
 \label{eq:tasks-wage-growth}
\end{equation}
If $p_{X,t}$ converges to a finite positive limit, average wage growth
converges to $\gamma$. Under competitive provision,
$w_t=A_t e^{\beta m_t}/B_t$, so average wage growth equals
$\gamma$ minus average efficiency growth whenever the latter has a limit.
In particular, bounded $B_t\to\overline B>0$ gives average wage growth $\gamma$.
\end{proposition}
\begin{proof}
Take logarithms of~\eqref{eq:tasks-wage-boundary}. Since $0<m_t<1$,
the last term divided by $t-t_0$ vanishes. Under competition substitute
$p_X=1/B$. These are average growth statements, not assertions about
instantaneous derivatives without additional regularity and convergence.
\end{proof}

For reference, define $E=e^{\beta(\sigma-1)m}$ and $S=1-m+D_\sigma E$.
Log-differentiating inverse demand and quantity with respect to $m$ gives
\begin{equation}
 e_X=-\frac{\dd\log P}{\dd\log X}
 =\frac{\alpha[\sigma\beta D_\sigma E/S+1/(1-m)]+\beta(1-m)/S}
 {D_\sigma'/D_\sigma+\sigma\beta+1/(1-m)}\longrightarrow\alpha.
 \label{eq:tasks-demand-elasticity}
\end{equation}
At $\sigma=1$, this formula is evaluated with $D_1=m$, $E=S=1$.
If a monopoly equilibrium has $X/H\to\infty$ and $B\to\overline B<\infty$,
its static supply condition implies
\[
 p_X\to\frac1{(1-\alpha)\overline B},\qquad
 \frac wA\to\frac{e^\beta}{(1-\alpha)\overline B}.
\]
Thus the original permanent wage-growth premium does not survive in this
benchmark, even though the aggregate capital-accumulation mechanism may survive.

The limiting capital return candidates are
\begin{align}
 \mathcal R^{\mathrm{mon}}(B)&=
 \alpha[B(1-\alpha)^2c_\sigma]^{(1-\alpha)/\alpha}-\delta,\\
 \mathcal R^{\mathrm{mc}}(B)&=
 \alpha[B(1-\alpha)c_\sigma]^{(1-\alpha)/\alpha}-\delta.
\end{align}
They follow from $X=BU$, the limiting markup, the AI revenue condition and
$\mathcal Z\sim c_\sigma X$. At a common target return the monopoly efficiency
requirement is $1/(1-\alpha)$ times the competitive requirement.
These are conditional limiting identities. They do not prove that either
trajectory exists or reaches that limit.

One implication is useful for the market comparison. Suppose both a competitive
trajectory with constant efficiency $B_0$ and a monopoly trajectory with
$B\to\overline B$ exist and satisfy $X/H\to\infty$, with the same task
technology and the same exogenous path of $A$. Then
\begin{equation}
 \frac{w_t^{\mathrm{mon}}}{w_t^{\mathrm{mc}}}
 \longrightarrow \frac{B_0}{(1-\alpha)\overline B},\qquad
 \mathcal R^{\mathrm{mon}}(\overline B)>\mathcal R^{\mathrm{mc}}(B_0)
 \quad\Longleftrightarrow\quad (1-\alpha)\overline B>B_0.
 \label{eq:tasks-wage-growth-ranking}
\end{equation}
Thus the larger limiting capital return goes with the lower limiting wage
relative to $A$ in this benchmark. Wage growth averages $\gamma$ under both
structures. These comparisons are conditional on the two trajectories and are
not welfare rankings or proof that such a pair exists from a common state.

\section{Proposed equilibrium and unfinished work}
I extend the original equilibrium definition, labelled
\texttt{def:rewrite-equilibrium} in its source.
Initial states are $K_0,B_0,A_0,N_0$ and population and labor-augmenting
technology retain their original laws. The household retains its original
preferences, budget, Euler condition and transversality requirement.

\begin{definition}[Task-automation equilibrium]
An equilibrium is an infinite-horizon trajectory of finite allocations and
prices such that the household maximizes utility, final-good firms maximize
profit using~\eqref{eq:tasks-aggregate}--\eqref{eq:tasks-technology}, and the
developer globally maximizes discounted net profit using the original research
technology. Labor, capital and AI-service markets clear, $H=AN$, $X=BU$,
the task allocation is efficient, and
\[
 C+\dot K+\delta K+U+M=Y,\qquad
 \dot B=\chi(BM)^\eta(1-B/\overline B).
\]
The household and developer transversality conditions and the relevant
corner optimality conditions must hold. Research is available from the initial
date. Neither $M>0$ nor an exogenous research-activation date is imposed.
\end{definition}

The developer faces $P(X;K,H)$ derived above, not the original fixed-weight
CES demand. At differentiability points the profit envelope still gives
$\pi_B=X/B^2$. Therefore the form of the interior research FOC and costate
can be retained, but their sufficiency must be proved again. The developer
continues to take aggregate capital and interest-rate paths as given.

Before numerical transitions I need to establish static monopoly uniqueness
or a justified selection, global research optimality, and dynamic equilibrium
existence with appropriate terminal conditions. I must also compare these
results with Korinek--Suh and Jones--Tonetti before making any novelty claim.
Distributions with unbounded relative difficulty or permanently human-only
tasks may change the wage bound. Automating production does not itself
endogenize automation of research.
\end{document}
